\section{Problems and certificates}\label{sec:model}

\subsection{The problem and its decomposition}
Let $X=\prod_{i=1}^n[L_i,U_i]$ be a compact box and consider
\begin{equation}\label{eq:model-problem}
 f^*=\min_{x\in X}F(x),\qquad
 F(x)=\sum_{t\in T}a_t(x_{V_t}).
\end{equation}
The set $T$ is the vertex set of a rooted tree, with root $r$ and bags
$V_t\subseteq\{1,\ldots,n\}$. The bags cover the coordinates. For each
coordinate $i$, the bags containing $i$ induce a connected subtree. This
is the running intersection property. A given factorization of $F$ can be
put in this form by assigning each factor to one bag containing its scope.
The analysis preserves this assignment; it does not merge bags or factors.

Write $\pi(t)$ for the parent of $t\ne r$, $\operatorname{ch}(t)$ for its
children, and $\sub(t)$ for its rooted subtree. Define
\[
 S_t=V_t\cap V_{\pi(t)},\quad R_t=V_t\setminus S_t\quad(t\ne r),
 \qquad S_r=\varnothing,\quad R_r=V_r.
\]
The relevant structural parameters are
\[
 N=|T|,\qquad p=\max_t|V_t|,\qquad
 q=\max_{t\ne r}|S_t|,\qquad
 k=\max_i|\{t:i\in V_t\}|.
\]
We take $q=0$ when $N=1$. In particular, $q\le p$; repeated or contained
bags are permitted. The width of this decomposition is $w=p-1$; the
graph's treewidth may be smaller. A bound in terms of $p$ alone does not
control $k$.

For a coordinate set $I$, $X_I$ denotes the projection of $X$ onto $I$.
For a box $B$, $B_I$ has the same meaning, and $\width(B)$ is its largest
side length. Fixed coordinates are substituted into the objective and
removed before constructing partitions. All remaining intervals have
positive length. A zero-dimensional projection has one point, one cell,
and width zero. If all coordinates are fixed, evaluation solves the problem.
Below $n\ge1$, $p\ge1$, and $s_0=\max_i(U_i-L_i)>0$.

Initially the functions $a_t$ need only be continuous. Their conditional
value functions are
\[
 \phi_t(s)=\min\left\{\sum_{u\in\sub(t)}a_u(x_{V_u}):
 x\in X_{\cup_{u\in\sub(t)}V_u},\ x_{S_t}=s\right\}.
\]
Running intersection and the product domain give the recursion
\begin{equation}\label{eq:model-value}
 \phi_t(s)=\min_{y\in X_{R_t}}
 \left[a_t(s,y)+\sum_{u\in\operatorname{ch}(t)}\phi_u((s,y)_{S_u})\right],
 \qquad \phi_r=f^*.
\end{equation}
Every minimum is attained. Induction in this recursion also shows that
$\phi_t$ is continuous: a continuous function minimized over a fixed compact
box has a continuous minimum. No differentiability of $\phi_t$ is assumed.

\subsection{Local models and growth}
For every bag box $B\subseteq X_{V_t}$, let $m_{t,B}$ be a supplied
continuous convex lower model of $a_t$ on $B$. The quantitative results use
the uniform contract
\begin{equation}\label{eq:model-contract}
 0\le a_t(v)-m_{t,B}(v)\le \frac A4\width(B)^2
 \qquad(v\in B),
\end{equation}
with $A\ge0$. This is a bag-level assumption. Factorwise quadratic-error
models imply it after summing their error bounds, but vertex exactness is
not required. The arithmetic implementation below uses affine Taylor
models that satisfy precisely this contract.

We also assume that $a_t$ is continuously differentiable on a neighborhood
of $X_{V_t}$ and
\begin{equation}\label{eq:model-smooth}
 \norm{\nabla a_t(v)-\nabla a_t(v')}\le M\norm{v-v'}
 \qquad(v,v'\in X_{V_t}),
\end{equation}
where $M\ge0$ and all unmarked norms are Euclidean. The main convergence
assumption is global quadratic growth about one minimizer:
\begin{equation}\label{eq:model-growth}
 F(x)-f^*\ge g\norm{x-x^*}^2\qquad(x\in X),\qquad g>0.
\end{equation}
It implies that $x^*$ is unique. The algorithm does not need $x^*$, and a
search variant does not need the numerical value of $g$. The minimizer may
lie on the boundary of $X$. Quadratic growth is an assumption on the whole
box; local positive curvature alone would not give the stated global
complexity bound without control of other nearly optimal regions.
We write $\log_+z=\max\{0,\log z\}$ for $z>0$.

\subsection{A finite lower certificate}
\begin{definition}\label{def:model-certificate}
A separator certificate consists of the following objects.
For each bag $t$, a finite partition $\mathcal L_t$ of $X_{V_t}$ into
closed boxes with disjoint interiors is specified. For each $t\ne r$,
a similar partition $\mathcal P_t$ of $X_{S_t}$ is specified, with affine
functions
\[
 l_{t,D}(s)=\lambda_{t,D}^{\mathsf T}s+\beta_{t,D}
 \qquad(D\in\mathcal P_t).
\]
For each $B\in\mathcal L_t$ and $u\in\operatorname{ch}(t)$, a continuous
convex function $\psi_{u,B}$ on $B_{S_u}$ satisfies
\begin{equation}\label{eq:model-child}
 \psi_{u,B}(s)\le l_{u,D'}(s)
 \quad\text{for every }D'\in\mathcal P_u, s\in B_{S_u}\cap D'.
\end{equation}
With $\operatorname{Rel}_{t,B}(v)=m_{t,B}(v)+
\sum_{u\in\operatorname{ch}(t)}\psi_{u,B}(v_{S_u})$, the local inequalities are
\begin{equation}\label{eq:model-local}
 \operatorname{Rel}_{t,B}(v)\ge l_{t,D}(v_{S_t})
 \quad(v\in B,\ v_{S_t}\in D),\qquad t\ne r, D\in\mathcal P_t.
\end{equation}
At the root they read $\operatorname{Rel}_{r,B}(v)\ge\LB$ for every
$B\in\mathcal L_r$ and $v\in B$.
The number of boxes and cells is
\[
 \mathcal S=\sum_t|\mathcal L_t|+\sum_{t\ne r}|\mathcal P_t|.
\]
A certificate together with a feasible value $\UB\ge f^*$ proves absolute
accuracy $\eps$ if $\UB-\LB\le\eps$.
\end{definition}

All nonempty intersections in this definition are included, even if they
are faces of smaller dimension. This convention avoids boundary choices
in the later configuration identity. The number $\mathcal S$ measures
partition size, not the number of local checks or the number of bits.
Representations of the models and local proof data must be counted
separately. A local check minimizes a continuous convex function over
$B\cap\{v:v_{S_t}\in D\}$, a compact box that may have fixed coordinates.

\begin{proposition}[Validity]\label{prop:model-validity}
Every certificate satisfies $l_{t,D}(s)\le\phi_t(s)$ for $s\in D$ and
$\LB\le f^*$.
\end{proposition}
\begin{proof}
Proceed from the leaves of the tree to the root. Fix $s\in D$ and choose
$v=(s,y)$ attaining \eqref{eq:model-value}. Some $B\in\mathcal L_t$
contains $v$. For each child $u$, choose a cell $D'\in\mathcal P_u$
containing $v_{S_u}$. The induction hypothesis and
\eqref{eq:model-child} imply
$\psi_{u,B}(v_{S_u})\le\phi_u(v_{S_u})$. Therefore
\[
 l_{t,D}(s)\le\operatorname{Rel}_{t,B}(v)
 \le a_t(v)+\sum_u\phi_u(v_{S_u})=\phi_t(s).
\]
The same argument at the root gives $\LB\le\phi_r=f^*$.
\end{proof}

\begin{proposition}[Cancellation along a consistent point]\label{prop:model-chain}
For any $x\in X$, choose a bag leaf $B_t$ containing $x_{V_t}$ at every
bag. Then
\begin{equation}\label{eq:model-chain}
 \LB\le\sum_t m_{t,B_t}(x_{V_t}).
\end{equation}
\end{proposition}
\begin{proof}
Choose any $D_t\in\mathcal P_t$ containing $x_{S_t}$. The local inequality
at $t$ and the child inequality at its parent apply to these choices,
because both projected boxes contain $x_{S_t}$. Substitute the local
inequality into the parent inequality, starting at the root. Every
nonroot affine minorant cancels once, leaving \eqref{eq:model-chain}.
\end{proof}

\subsection{A dynamic program for fixed slopes}
The construction uses one slope $\lambda_t$ for all cells of a separator.
For $B\in\mathcal L_t$, set
\begin{equation}\label{eq:model-fixed-child}
 \psi_{u,B}(s)=\lambda_u^{\mathsf T}s+
 \min_{D'\in\mathcal P_u:\ D'\cap B_{S_u}\ne\varnothing}\beta_{u,D'}.
\end{equation}
This affine function satisfies \eqref{eq:model-child}. In postorder,
compute the largest permissible intercepts
\begin{equation}\label{eq:model-intercept}
 \beta_{t,D}=\min_{B\in\mathcal L_t:\ B_{S_t}\cap D\ne\varnothing}
 \min_{v\in B:\ v_{S_t}\in D}
 \left[m_{t,B}(v)+\sum_{u\in\operatorname{ch}(t)}\psi_{u,B}(v_{S_u})
       -\lambda_t^{\mathsf T}v_{S_t}\right],
\end{equation}
and at the root compute $\LB=\min_{B\in\mathcal L_r}
\min_{v\in B}\operatorname{Rel}_{r,B}(v)$. These are convex minimization
problems. They have minimizers because their domains are nonempty compact
boxes and their objectives are continuous. The calculation produces a
certificate for any choice of slopes and partitions.

A \emph{configuration} chooses $B_t\in\mathcal L_t$ and $z^t\in B_t$
for every bag, and $D_t\in\mathcal P_t$ for every nonroot bag, with
\[
 z^t_{S_t}\in D_t,\qquad D_t\cap(B_{\pi(t)})_{S_t}\ne\varnothing.
\]
It need not equate the two copies of a separator coordinate.
\begin{proposition}[Unfolding the dynamic program]\label{prop:model-dp}
The fixed-slope program satisfies
\begin{equation}\label{eq:model-config}
 \LB=\min_{\text{configurations}}\Phi,\qquad
 \Phi=\sum_t m_{t,B_t}(z^t)+
 \sum_{t\ne r}\lambda_t^{\mathsf T}
       (z^{\pi(t)}_{S_t}-z^t_{S_t}).
\end{equation}
A minimizing configuration can be recovered from stored minimizing
choices and local minimizers.
\end{proposition}
\begin{proof}
In \eqref{eq:model-intercept}, substitute \eqref{eq:model-fixed-child}.
For each child the minimizing cell is constrained only to meet the
projection of $B$, independently of the parent's point $v$. Substituting
the child's intercept introduces its own leaf and point, constrained by
that cell. Continue down the tree. A sum of independent minima equals
the minimum over their joint choices. Each nonroot slope occurs once
at the parent point and once with a minus sign at its own point.
The resulting expression is exactly \eqref{eq:model-config}. All choices
are among finitely many compact domains, so the minimum is attained.
Storing the choices during the postorder pass and following them from
the root recovers a minimizing configuration.
\end{proof}

\subsection{Slopes computed at a current center}
Let $\topbag(i)$ be the bag closest to the root among those containing
$i$. Running intersection makes it unique. A configuration defines a
consistent point by $x_i=z_i^{\topbag(i)}$. This point belongs to $X$.
For any center $c\in X$, define
\begin{equation}\label{eq:model-slope}
 \lambda_{t,i}(c)=
 \sum_{u\in\sub(t):\ i\in V_u}\partial_i a_u(c_{V_u})
 \qquad(i\in S_t).
\end{equation}
These are subtree sums of ordinary bag gradients, not derivatives of
conditional value functions. They can be computed by a postorder pass
over the coordinate occurrences.

\begin{lemma}[First-order cancellation]\label{lem:model-telescope}
For every center and every configuration,
\begin{equation}\label{eq:model-telescope}
 \sum_t\nabla a_t(c_{V_t})^{\mathsf T}(z^t-x_{V_t})
 =\sum_{t\ne r}\lambda_t(c)^{\mathsf T}
       (z^t_{S_t}-z^{\pi(t)}_{S_t}).
\end{equation}
Consequently, if $e_t(v)=a_t(v)-m_{t,B_t}(v)$ and the slopes are
\eqref{eq:model-slope}, then
\begin{align}\label{eq:model-error-identity}
 \Phi=F(x)+\sum_t\bigl[
 a_t(z^t)-a_t(x_{V_t})
 -\nabla a_t(c_{V_t})^{\mathsf T}(z^t-x_{V_t})\bigr]
 -\sum_t e_t(z^t).
\end{align}
\end{lemma}
\begin{proof}
Fix $i$. On its occurrence subtree, $z_i^t-x_i$ is the sum of the
successive differences $z_i^u-z_i^{\pi(u)}$ on the path from
$\topbag(i)$ to $t$. Multiply by $\partial_i a_t(c_{V_t})$ and sum over
the occurrences $t$. The coefficient of a given edge difference is the
sum of these derivatives over the occurrences below that edge, which
is $\lambda_{u,i}(c)$. Summing over coordinates proves
\eqref{eq:model-telescope}. Substitution into \eqref{eq:model-config}
gives \eqref{eq:model-error-identity}.
\end{proof}

The cancellation holds at every center in the box. It uses neither
stationarity nor knowledge of $x^*$. This is the identity that allows a
certificate to guide the choice of the next center.
